\section{BCA-Style Accounting Logic}

The accounting exercise follows the spirit of Business Cycle Accounting. The
goal is not to claim that a wedge is itself a deep friction. The goal is to
measure which equilibrium condition must be distorted in order for the simple
model to reproduce the observed dynamics.

\subsection{Step 1: Infer Wedges}

Given observed or model-implied series for
\(\{x_t,\pi_t,i_t,r^{n,1}_t\}\) and expectations
\(\{\E_t x_{t+1},\E_t \pi_{t+1}\}\), infer the wedges as residuals:
\begin{align}
    \widehat{\omega}^x_t
    &=
    x_t
    -
    \E_t x_{t+1}
    +
    \frac{1}{\gamma}
    \left(
        i_t-\E_t\pi_{t+1}-r^{n,1}_t
    \right),
    \label{eq:infer_x_wedge}
    \\
    \widehat{\omega}^\pi_t
    &=
    \pi_t
    -
    \beta \E_t\pi_{t+1}
    -
    \kappa x_t,
    \label{eq:infer_pi_wedge}
    \\
    \widehat{\omega}^i_t
    &=
    i_t
    -
    \rho_i i_{t-1}
    -
    (1-\rho_i)(\phi_\pi\pi_t+\phi_x x_t).
    \label{eq:infer_i_wedge}
\end{align}

\subsection{Step 2: Test the Single-Uncertainty-Source Restriction}

Estimate or inspect the relation
\begin{equation}
    \widehat{\bs{\omega}}_t
    =
    \bs{\lambda} q_t
    +
    \bs{u}_t.
    \label{eq:estimate_wedge_loadings}
\end{equation}
The single-source TFP uncertainty hypothesis implies that most systematic
variation in the wedge vector should be explained by the scalar variance state
\(q_t\). In this minimal model, a large and systematic \(\widehat{\omega}^i_t\)
would be evidence against the clean baseline, because monetary policy is not a
primitive uncertainty source.

\subsection{Step 3: Counterfactual Wedge Experiments}

Once the wedges are inferred, run counterfactual paths by opening one wedge at a
time:
\begin{enumerate}[label=(\alph*)]
    \item Demand-risk counterfactual: set
    \(\omega^x_t=\widehat{\omega}^x_t\) and
    \(\omega^\pi_t=\omega^i_t=0\).
    \item Markup-risk counterfactual: set
    \(\omega^\pi_t=\widehat{\omega}^\pi_t\) and
    \(\omega^x_t=\omega^i_t=0\).
    \item Policy-residual counterfactual: set
    \(\omega^i_t=\widehat{\omega}^i_t\) and
    \(\omega^x_t=\omega^\pi_t=0\).
\end{enumerate}
The contribution of each wedge is evaluated by how closely the counterfactual
paths reproduce the observed paths of \(x_t\), \(\pi_t\), and \(i_t\).

\subsection{Connection to the Later SOE Nonlinear BCA}

The later small-open-economy framework will contain more wedges: efficiency,
labor, investment, trade, and UIP wedges. This note is a stripped-down
laboratory. Its value is to establish the logic:
\[
    \text{second-moment primitive}
    \rightarrow
    \text{equation-specific risk wedge}
    \rightarrow
    \text{counterfactual accounting}.
\]
Once this logic is clear, the same structure can be expanded to multiple
external uncertainty factors and multiple small-open-economy wedges.
